% =====================================================================
%  Introduction générale
% =====================================================================
\chapter*{Introduction générale}
\addcontentsline{toc}{chapter}{Introduction générale}
\markboth{Introduction générale}{}
\label{chap:introduction}

\section*{Un secteur sous double contrainte}

L'industrie pharmaceutique se transforme sous deux pressions simultanées. D'un côté,
la pression économique : les coûts de développement d'un médicament augmentent, les
marges se resserrent, et l'optimisation des opérations devient un levier de
compétitivité. De l'autre, la pression réglementaire : chaque processus touchant au
produit, au lot ou au dossier patient relève d'un cadre GxP qui impose traçabilité,
validation informatique et justification des décisions.

L'intelligence artificielle est présentée comme une réponse à la première pression.
Elle se heurte pourtant frontalement à la seconde. Un modèle dont on ne peut pas
expliquer la sortie est difficilement défendable devant un inspecteur. Une
automatisation qui décide seule sur un périmètre GxP est encore plus difficile à
justifier.

\section*{Le problème observé}

Dans les organisations pharmaceutiques, l'adoption de l'IA échoue rarement par manque
d'idées. Elle échoue plus souvent en amont du code, pour des raisons qui n'ont rien de
technique :

\begin{itemize}
  \item \textbf{les besoins remontent de départements cloisonnés.} Le Commercial, la
    Production, la Qualité et les Affaires réglementaires expriment leurs irritants
    chacun de leur côté, souvent dans des formats différents, parfois sur le même
    processus sans le savoir ;
  \item \textbf{la qualification technique arrive trop tard.} On décide de « faire de
    l'IA » avant de vérifier si les données existent, si elles sont exploitables et si
    une vérité terrain est disponible. Le projet s'arrête alors au premier atelier
    technique ;
  \item \textbf{les arbitrages manquent de critères partagés.} Sans grille commune,
    la priorisation se joue à la conviction du sponsor le plus influent plutôt qu'à
    l'impact attendu ;
  \item \textbf{la valeur n'est pas mesurée.} Le business case sert à obtenir le budget,
    puis disparaît. Personne ne compare la valeur promise à la valeur réalisée ;
  \item \textbf{la traçabilité fait défaut.} Six mois plus tard, plus personne ne sait
    pourquoi telle solution a été retenue plutôt qu'une autre.
\end{itemize}

Ces cinq difficultés partagent une même cause : l'absence d'un référentiel unique où
le diagnostic, la qualification, l'arbitrage et le suivi vivent ensemble, avec leurs
justifications.

\section*{La réponse apportée}

\appname{} est une plateforme SaaS multi-entreprises qui couvre le parcours complet,
de l'état des lieux à la valeur captée. Elle s'organise autour de sept temps
(figure~\ref{fig:parcours}).

\begin{figure}[H]
\centering
\begin{tikzpicture}[node distance=0.42cm and 0.5cm]
  \node[blocvert] (diag) {\textbf{Diagnostic}\\\scriptsize AS-IS, maturité};
  \node[blocvert, right=of diag] (qual) {\textbf{Qualification}\\\scriptsize type de solution};
  \node[blocvert, right=of qual] (read) {\textbf{Readiness}\\\scriptsize score sur 100};
  \node[blocvert, below=0.9cm of read] (prio) {\textbf{Priorisation}\\\scriptsize score, rang, vague};
  \node[blocvert, left=of prio] (fiche) {\textbf{Fiche projet}\\\scriptsize versions, décision};
  \node[blocvert, left=of fiche] (cycle) {\textbf{Cycle de vie}\\\scriptsize phases, gates};
  \node[blocsombre, below=0.9cm of fiche] (roi) {\textbf{Valeur \& ROI}\\\scriptsize TCO, gains, payback};

  \draw[flechevert] (diag) -- (qual);
  \draw[flechevert] (qual) -- (read);
  \draw[flechevert] (read) -- (prio);
  \draw[flechevert] (prio) -- (fiche);
  \draw[flechevert] (fiche) -- (cycle);
  \draw[flechevert] (cycle.south) |- (roi.west);
  \draw[fleche, dashed] (roi.east) -| node[pos=0.25, below, font=\scriptsize, text=medicaMuted]
    {rétroaction : la valeur mesurée corrige les hypothèses} (diag.south);
\end{tikzpicture}
\caption{Le parcours couvert par la plateforme, du diagnostic à la mesure de la valeur}
\label{fig:parcours}
\end{figure}

Chacun de ces temps produit une décision, et chaque décision reste rattachée à ce qui
l'a motivée. Le readiness renvoie la note de chaque dimension avec sa preuve. La
recommandation de solution expose les règles satisfaites et les alternatives écartées.
Le score de priorité détaille la contribution de ses six critères. Le business case
publie ses coûts et ses gains ligne à ligne.

\section*{Ce que ce rapport établit}

Ce document décrit ce qui a été \emph{réellement} développé. Lorsque l'implémentation
s'écarte du cahier des charges initial, l'écart est signalé explicitement dans un
encadré dédié, avec ce qui était demandé, ce qui a été réalisé et ce qui reste à faire.

Un point mérite d'être posé dès l'introduction, car il conditionne la lecture de tout
le rapport. Le nom du produit contient « IA », et la plateforme recommande des types
de solutions d'intelligence artificielle. Elle ne \emph{contient} pourtant aucun modèle
d'apprentissage automatique : ni modèle entraîné localement, ni appel à un fournisseur
externe. Sa capacité de recommandation repose sur un moteur de règles pondérées,
déterministe et explicable. Ce choix est argumenté à la
section~\ref{sec:pas-dia}~; le confondre avec de l'apprentissage automatique serait une
erreur d'analyse.

\section*{Organisation du document}

Le rapport suit la progression du projet.

\begin{itemize}
  \item Le \textbf{chapitre~\ref{chap:analyse}} établit le contexte, la problématique,
    les objectifs, le cahier des charges et l'état réel de chacun des douze modules
    fonctionnels.
  \item Le \textbf{chapitre~\ref{chap:conception}} présente les acteurs, le modèle de
    droits, les cas d'utilisation, le modèle de données et l'architecture technique.
  \item Le \textbf{chapitre~\ref{chap:realisation}} détaille les technologies employées,
    les algorithmes de readiness, de priorisation et de valorisation, le moteur de
    qualification et le dispositif de sécurité.
  \item Le \textbf{chapitre~\ref{chap:interfaces}} parcourt les quatorze écrans de
    l'application. Chaque écran est accompagné d'une proposition de présentation orale
    directement réutilisable en soutenance.
  \item Le \textbf{chapitre~\ref{chap:tests}} rend compte de la stratégie de validation
    et des résultats obtenus.
  \item Le \textbf{chapitre~\ref{chap:conclusion}} dresse le bilan, énonce les limites
    de la version actuelle et ouvre sur les évolutions envisageables.
\end{itemize}

L'annexe~\ref{ann:soutenance} propose un guide de démonstration minuté et une série de
questions susceptibles d'être posées par le jury, avec des réponses appuyées sur le
code réel.
